%
% sn-konstruktion.tex -- geometrische Definition der Jacobischen
% elliptischen Funktionen (nach Mueller, Spezielle Funktionen 2022,
% Abb. 11.4/11.5)
%
% Links: Einheitskreis, Bogenlaenge u = 1.0, Punkt (cos u, sin u).
% Rechts: Ellipse x^2/a^2 + y^2 = 1 mit a = 1/sqrt(1-k^2), k = 0.8,
% Hilfskreis (gestrichelt). Q = (cos phi, sin phi) auf dem Kreis,
% P = (a cos phi, sin phi) auf der Ellipse, phi = am(u,k).
% sn(u,k), cn(u,k) und r = a dn(u,k) erscheinen als Strecken.
%
% Zahlenwerte: phi = am(1.0, 0.8) numerisch bestimmt per RK4 auf
%   dphi/du = sqrt(1 - k^2 sin^2(phi)), phi(0) = 0, Schrittweite 1e-5,
%   Kontrolle per Simpson-Quadratur von
%   u(phi) = int_0^phi dtheta/sqrt(1 - 0.64 sin^2(theta)) = 1.00000.
% Ergebnis:
%   phi     = 0.913564 (52.343 Grad)
%   cos phi = 0.610928  (= cn(u,k))
%   sin phi = 0.791686  (= sn(u,k))
%   dn(u,k) = 0.773866,  r = a*dn = 1.289777
% Linkes Panel: u = 1.0 rad = 57.296 Grad,
%   Punkt = (0.540302, 0.841471).
%
\documentclass[tikz]{standalone}
\usepackage{amsmath}
\usepackage{times}
\usepackage{txfonts}
\definecolor{darkred}{rgb}{0.8,0,0}
\definecolor{darkblue}{rgb}{0,0,0.6}
\begin{document}
\def\skala{2.0}
\begin{tikzpicture}[>=latex, thick, scale=\skala]

% numerische Werte (siehe Kommentar oben)
\def\a{1.66667}        % grosse Halbachse 1/sqrt(1-k^2), k = 0.8
\def\cq{0.610928}      % cos(phi), phi = am(1.0, 0.8)
\def\sq{0.791686}      % sin(phi)
\def\px{1.018213}      % a*cos(phi), x-Koordinate von P

% ---------- linkes Panel: Einheitskreis ----------
% helle Achsen
\draw[->, black!45] (-1.35,0) -- (1.55,0) node[anchor=north east] {$x$};
\draw[->, black!45] (0,-1.35) -- (0,1.55) node[anchor=south east] {$y$};

% Einheitskreis
\draw (0,0) circle (1);

% Radius zum Punkt, duenn
\draw[thin] (0,0) -- (0.540302,0.841471);

% Kreisbogen von (1,0) bis zum Punkt = Bogenlaenge u
\draw[darkred, line width=1.6pt]
	(1,0) arc[start angle=0, end angle=57.296, radius=1];
\node[darkred] at (28.65:1.22) {$u$};

% Punkt (cos u, sin u)
\fill[darkred] (0.540302,0.841471) circle (0.035);
\node[darkred, anchor=south west] at (0.52,0.86) {$(\cos u, \sin u)$};

% ---------- rechtes Panel: Ellipse mit Hilfskreis ----------
\begin{scope}[xshift=4.3cm]

% Achsen
\draw[->, black!45] (-1.9,0) -- (2.1,0) node[anchor=north east] {$x$};
\draw[->, black!45] (0,-1.35) -- (0,1.6) node[anchor=south east] {$y$};

% Hilfskreis (gestrichelt)
\draw[dashed, thin] (0,0) circle (1);

% Ellipse x^2/a^2 + y^2 = 1
\draw (0,0) ellipse ({\a} and 1);

% Halbachsen markieren
\draw[thin] (0,0) -- (\a,0) node[midway, below] {$a$};
\fill (\a,0) circle (0.03);
\draw[thin] (0,0) -- (0,1) node[midway, left] {$1$};
\fill (0,1) circle (0.03);

% Linie vom Ursprung zu Q und Winkel phi
\draw[thin] (0,0) -- (\cq,\sq);
\draw[thin] (0.32,0) arc[start angle=0, end angle=52.343, radius=0.32];
\node at (26.17:0.42) {$\varphi$};

% sn(u,k): vertikale gestrichelte Linie von P zur x-Achse
\draw[dashed, thin] (\px,0) -- (\px,\sq)
	node[midway, sloped, below] {$\operatorname{sn}(u,k)$};

% cn(u,k): horizontale gestrichelte Linie von der y-Achse zu Q
\draw[dashed, thin] (0,\sq) -- (\cq,\sq)
	node[pos=0.3, below, fill=white, inner sep=1.5pt] {$\operatorname{cn}(u,k)$};

% gestrichelte horizontale Verbindung Q -- P
\draw[dashed, thin] (\cq,\sq) -- (\px,\sq);

% Radius vom Ursprung zu P
\draw[darkred] (0,0) -- (\px,\sq)
	node[pos=0.62, below=3pt, darkred] {$r$};

% Bedeutung des Radius
\node[darkred] at (0.85,1.32) {$r = a\,\operatorname{dn}(u,k)$};

% Punkt Q auf dem Hilfskreis
\fill (\cq,\sq) circle (0.028);
\node[anchor=north, yshift=1.5pt] at (\cq,\sq) {$Q$};

% Punkt P auf der Ellipse
\fill[darkred] (\px,\sq) circle (0.035);
\node[darkred, anchor=south west] at (\px,\sq) {$P$};

\end{scope}

\end{tikzpicture}
\end{document}
